\medskip










Une chaîne de fabrication produit des pièces mécaniques. On estime que 4\,\% des pièces produites par cette chaîne sont défectueuses.

On choisit au hasard $n$ pièces produites par la chaîne de fabrication. 

Le nombre de pièces produites est suffisamment grand pour que ce choix puisse être assimilé à un tirage avec remise. 

On note $X$ la variable aléatoire égale au nombre de pièces défectueuses tirées.

Dans les trois questions suivantes, on prend $n = 50$.

\medskip

\begin{enumerate}
\item 
On a $p(X\geqslant 1) = 1 - p(X = 0) = 1 - 0,04^0 \times 0,96^{50} \approx 1 - \np{0,1299} \approx \np{0,8701}$ soit 0,870 au millième près : réponse \textbf{b.~~}






\item La probabilité $p(3 < X \leqslant 7)$ est égale à :
$p(X \leqslant 7) - p(x \leqslant 3) \approx \np{0,9992} - \np{0,8609}$ soit \np{0,1383}. Réponse \textbf{b.~~}








\item Quel est le plus petit entier naturel $k$ tel que la probabilité de tirer au plus $k$ pièces défectueuses soit supérieure ou égale à $95\,\%$ ?







La calculatrice donne $k = 4$, réponse \textbf{c.}

Dans les questions suivantes, $n$ ne vaut plus nécessairement $50$.

\begin{enumerate}[resume]
\item Quelle est la probabilité de ne tirer que des pièces défectueuses ?






Réponse \textbf{a.}
\item On considère la fonction Python ci-dessous. Que renvoie-t-elle ?
















La fonction donne le plus petit naturel $n$ tel que $1 - 0,96^n \geqslant $ x.

Réponse \textbf{a.}
	\end{enumerate}
\end{enumerate}
